%%% Header
\documentclass[11pt,a4paper,titlepage,
chapterprefix,headsepline,parskip,pdftex,
,pointlessnumbers,bibtotoc]{article}

%%% paragraph for lower depth
\makeatletter %% remove @ meaning
\renewcommand{\paragraph}{\@startsection
   {paragraph} % name
   {4} % depth
   {0mm} % indention
   {-\baselineskip} % before
   {0.1\baselineskip} % after
   {\normalfont\normalsize\bfseries}} % stil
\makeatother %% add @ meaning

\makeatletter %% remove @ meaning
\renewcommand{\subparagraph}{\@startsection
   {subparagraph} % name
   {5} % depth
   {0mm} % indention
   {-\baselineskip} % before
   {0.1\baselineskip} % after
   {\normalfont\normalsize\bfseries}} % stil
\makeatother %% add @ meaning

\usepackage{setspace}
\onehalfspacing

\usepackage[pdftex]{graphicx}

% for colours
\usepackage[pdftex]{color}

% Configure defintions
\usepackage{amsthm}
\newtheorem{mydef}{Definition}

\usepackage{tikz}
\usetikzlibrary{shapes,arrows}
% \usepackage{subfigure} // cannot be used with subcaption
\usepackage{caption}
\usepackage{subcaption}

\usepackage[colorlinks=true,
    linkcolor=black,
    citecolor=black,
    pagecolor=black,
    urlcolor=black,
    breaklinks=true,
    bookmarksnumbered=true,
    hypertexnames=false,
    pdfpagemode=UseOutlines,
    pdfview=FitH,
    plainpages=false,
    pdfpagelabels,
    bookmarks=true,
    linktocpage=true]{hyperref}

\hypersetup{pdfauthor={Markus Hiesmair},
    pdftitle={Austrian Parliament Analyzer},
    pdfsubject={Bachelor's Thesis},
    pdfkeywords={},
    pdfcreator={pdfLaTeX with hyperref (\today})}

%%% Source-Code
\usepackage{listings}
\lstset{language=Java}

\lstset{% general command to set parameter(s)
basicstyle=\small, % print whole listing small
tabsize=2, %
keywordstyle=\color[rgb]{0.00,0.00,0.50}{}\bfseries,
%identifierstyle=, % nothing happens
commentstyle=\color[rgb]{0.00,0.50,0.25}{},
%stringstyle=\ttfamily, % typewriter type for strings
%showstringspaces=false} % no special string spaces
numbers=left, numberstyle=\tiny, numbersep=5pt}

\colorlet{punct}{red!60!black}
\definecolor{background}{HTML}{EEEEEE}
\definecolor{delim}{RGB}{20,105,176}
\colorlet{numb}{magenta!60!black}

% Custom definition for JSON
\lstdefinelanguage{json}{
    basicstyle=\normalfont\ttfamily,
    numbers=left,
    numberstyle=\scriptsize,
    stepnumber=1,
    numbersep=8pt,
    showstringspaces=false,
    breaklines=true,
    frame=lines,
    backgroundcolor=\color{background},
    literate=
     *{0}{{{\color{numb}0}}}{1}
      {1}{{{\color{numb}1}}}{1}
      {2}{{{\color{numb}2}}}{1}
      {3}{{{\color{numb}3}}}{1}
      {4}{{{\color{numb}4}}}{1}
      {5}{{{\color{numb}5}}}{1}
      {6}{{{\color{numb}6}}}{1}
      {7}{{{\color{numb}7}}}{1}
      {8}{{{\color{numb}8}}}{1}
      {9}{{{\color{numb}9}}}{1}
      {:}{{{\color{punct}{:}}}}{1}
      {,}{{{\color{punct}{,}}}}{1}
      {\{}{{{\color{delim}{\{}}}}{1}
      {\}}{{{\color{delim}{\}}}}}{1}
      {[}{{{\color{delim}{[}}}}{1}
      {]}{{{\color{delim}{]}}}}{1},
}

\lstset{
    basicstyle=\normalfont\ttfamily,
    numbers=left,
    numberstyle=\scriptsize,
    stepnumber=1,
    numbersep=8pt,
    showstringspaces=false,
    breaklines=true,
    frame=lines,
    backgroundcolor=\color{background},
}


% Figure
\newcommand{\cffigure}[1]{\hyperref[#1]
{cf. \ref*{#1}~\nameref{#1}}}


%%% continous footnote
\newcounter{cfootnotecounter}
\newcommand{\cfootnote}[1]{\stepcounter{cfootnotecounter}
\footnote[\value{cfootnotecounter}]{#1}}

\usepackage{footnote}
\makesavenoteenv{tabular}
\makesavenoteenv{table}

\flushbottom

% change page settings
\setlength{\hoffset}{0mm} \setlength{\voffset}{0mm}
\setlength{\evensidemargin}{14.6mm}
\setlength{\oddsidemargin}{14.6mm} \setlength{\topmargin}{-20mm}
\setlength{\headheight}{15mm} \setlength{\headsep}{9mm}
\setlength{\textheight}{242mm} \setlength{\textwidth}{145mm}
\setlength{\footskip}{10mm}


%%% seperation of float-environment
\setlength{\textfloatsep}{25pt plus5pt minus5pt}
\setlength{\intextsep}{25pt plus5pt minus5pt}

%%% Gliederungs-Nummern in den Rand schreiben
%\renewcommand*{\othersectionlevelsformat}[1]{%
%\llap{\csname the#1\endcsname\autodot\enskip}}

%%% only chapter name in headline
%\renewcommand*{\chaptermarkformat}{}

%%% format of the chapter looks
%\setkomafont{chapter}{\Huge}
%\renewcommand*{\chapterformat}{\LARGE{\chapappifchapterprefix{\ }\thechapter\autodot\enskip}}

%%% headline
\usepackage[automark]{scrpage2}

\clearscrheadings \clearscrplain \clearscrheadfoot
\pagestyle{scrheadings}
\ohead{\pagemark}
\ihead{\headmark}
\cfoot{}

%%% design of chapter pages
%\renewcommand*{\chapterpagestyle}{scrheadings}

%% indices from TOC and numbering depth
\setcounter{tocdepth}{\subsubsectionlevel}
\setcounter{secnumdepth}{\paragraphlevel}

%%% Array for tables
\usepackage{array}

%%% fonts
%\addtokomafont{chapter}{\sffamily}
%\addtokomafont{sectioning}{\rmfamily}

% language
\usepackage[english]{babel}
% inputencoding
\usepackage[utf8]{inputenc}
% fontencoding
\usepackage[T1]{fontenc}
\usepackage{ae}

% URLs
\usepackage{url}

%%% orphan  und widow
\clubpenalty = 10000
\widowpenalty = 10000 \displaywidowpenalty = 10000


%%% Einbinden von kompletten PDF-Seiten
\usepackage{pdfpages}

\usepackage{subcaption}
\expandafter\def\csname ver@subfig.sty\endcsname{}

\usepackage{ifluatex}
\ifluatex
  \usepackage{pdftexcmds}
  \makeatletter
  \let\pdfstrcmp\pdf@strcmp
  \let\pdffilemoddate\pdf@filemoddate
  \makeatother
\fi

\usepackage{svg}
\usepackage{amsmath}
\usepackage{xr}
\usepackage{epstopdf}
\usepackage{wrapfig}

% comment the following line for final submission
%\newcommand{\manuscriptmode}[0]{draft}

\newcommand{\thesistitle}[0]{Making Computers Understand Coalition and Opposition in Parliamentary Democracy}

%%% Custom commands
% provides support for a \todo command
% This file contains custom commands - makes the rest of the structure updateable

% A ToDo Command that is used only during drafts
\newcommand{\todo}[1]{\textcolor{red}{\textbf{To-Do:} #1}}

% Common picture
\newcommand{\cpic}[3]{\begin{figure}[!t]\centering
\includegraphics[width=3.5in]{#1}
\caption{#2} \label{#3}\end{figure}}

\newcommand{\cpicw}[4]{\begin{figure}[!t]\centering
\includegraphics[width=#4]{#1}
\caption{#2} \label{#3}\end{figure}}

% A command that highlights java classes
\newcommand{\inlinecode}[1]{\emph{#1}}





%%% hyphenation
%%% Hyphenation

\tolerance=2000                 % against "overfull hbox"
\emergencystretch 20pt          % better than sloppy"


\hyphenation{hy-phen-ate}


\begin{document}

%%% cover page
\pagenumbering{roman}
\includepdf{structure/cover/cover.pdf}

%%% structure parts
\pagenumbering{Roman}
%\ifx\manuscriptmode\undefined{}
	%%\null\vspace{\stretch{1}}
%\begin{flushright}
%To Sonja the love of my life!
%\end{flushright}
%\vspace{\stretch{2}}\null

%\fi
%\chapter*{Affidavit}
\markright{Affidavit}

I hereby declare that the following dissertation "\thesistitle{}" has been written
only by the undersigned and without any assistance from third parties.

Furthermore, I confirm that no sources have been used in the preparation of this thesis other than those indicated in the thesis itself. This printed thesis is identical with the electronic version submitted.

Linz, on \today

\hfill Markus Hiesmair

%\chapter*{Acknowledgment}
\markright{Acknowledgment}

\ifx\manuscriptmode\undefined{}

I would like to thank my supervisor Assoc. Prof. Dr. Karin Anna Hummel for her support throughout the whole work on this thesis. She always provided me useful feedback and used her professional skills and knowledge to give me hints and advices when I needed it. 

\fi

%%\chapter*{Summary}
%\markright{Summary}

%Summary \ldots

\chapter*{Abstract}
\markright{Abstract}

Traffic congestion in urban areas is getting a big problem nowadays. Not only the increased traffic contributes to traffic jams, but also drivers who are looking for parking spaces. When driving around a long time in the city to find vacant parking spaces, they not only congest the streets, but also emit a lot of green house gases.
% which are bad for the environment.
% (for studies of parking space induced traffic see \cite{Nawaz:2013:PSB:2500423.2500438, TexasMobilityReport}). 
To mitigate these problems, it would be helpful to increase the transparency of the current parking space availability situation to help drivers navigate directly to vacant parking spaces close to their end destination. 

This thesis tackles the task of creating parking space availability maps for cities to make it easier for drivers to find vacant parking spaces. A drive-by approach is chosen to detect vacant and occupied parking spaces. 
Using our approach, several vehicles are driving through the city and continuously collecting sensor data. All vehicles measure the distance to the right side of the road as well as their position using GPS while driving. Using the distance information, our prototype should be able to decide whether a parking space is occupied or vacant. Furthermore, distracting situations which prevent a successful detection (e.g. an overtaking situation) should be detected as well. A prototype vehicle has been equipped with the necessary sensors. Real world sensor data has been collected during 32 test drives. This data is used in the evaluation process to identify the best performing systems.

As the traffic situations in our test drives can be fairly complicated, it is non-trivial to find a rule set which can detect the different situations. Thus, machine learning techniques should be tested on their performance. Using a manually tagged dataset containing nine computed features derived from the sensor data, several traditional machine learning algorithms and deep learning approaches are evaluated on their classification performance. Furthermore, several techniques to further improve the classification performance are tested, such as over- and under-sampling. 
Finally, a custom two-staged classification process is introduced to include information of the surroundings of the parking situation to classify. This approaches reaches the best classification results. In total, 96.52\% of the situations are correctly classified and 93.81\% of the parking cars are detected.



%%%% Diverse Verzeichnisse

\setcounter{secnumdepth}{3}
\setcounter{tocdepth}{3}
\tableofcontents
%\todo{For now the TOC depth is 5 but will be reduced later}


\chapter*{Abbreviations}
\markright{Abbreviations}

\begin{description}
\setlength{\itemsep}{-11pt}
\setlength{\leftmargin}{900pt}

%C
\item[CNN] Convolutional Neural Network

%D
\item[DL] Deep Learning

%G
\item[GNNS] Global Navigation Satellite System
\item[GPIO] General Purpose Input Output
\item[GPS] Global Positioning System

%I
\item[I2C] Inter Integrated Circuit

%L
\item[LIDAR] Light detection and ranging

\end{description}


\listoffigures

\listoftables


%%% content
\newpage
\pagenumbering{arabic}

\section{Initial Situation and Motivation}
Currently the road side parking situation in most cities is rather untransparent. Except from parking garages and the like information about the availability of parking spaces is rarely available. However, finding parking spaces in urban areas can be a really difficult, frustrating and time consuming task for drivers. Furthermore, information about free parking spaces can help to reduce traffic by a tremendous amount. Studies have shown that in urban areas about 30\% of traffic congestion is created by drivers looking for free parking spaces \cite{Nawaz:2013:PSB:2500423.2500438} and that in 2007 a loss of about \$78 billion U.S. dollars was created by the use of about 2.9 billion gallons of gasoline alone in the USA \cite{TexasMobilityReport}. Furthermore, about 4.9 billion hours were wasted by drivers while looking for parking spaces during that year. 

\section{Problem Definition}

Detection of road side parking spaces and their states is a challenging task. Of course an obvious approach to the problem would be to put sensors to every parking space in the city, which check, if the corresponding parking space is occupied or vacant. This, however, has the drawback to be very expensive as, for big cities, thousands of sensors would have to be bought, installed and maintained. Furthermore, because the parking situation does not change often, the high frequency of sensing with such a system would be rather inefficient.

Another promising option to sense a city's parking situation is the use of mobile sensors instead of static ones. This approach has the advantage to be usually more cost effective and can provide sufficient accuracy.
There has already been done some research on cars which analyze parking availability while they drive through the city. For instance, Mathur et al. \cite{Mathur:2010:PDS:1814433.1814448} developed a system which uses distance information from the vehicle to the right side of the road to reason about parking spaces. Another approach, which was done by Zhou et al. \cite{Zhou_2d_range} looks for car bumper shaped signal parts in the distance measurements to identify parking cars.

However, all of the mentioned mobile sensing approaches only work in single lane scenarios. Multi lane streets bring much more complexity in the recognition of parking spaces. There are many special cases which have to be addressed to work properly on multi lane roads. For example, the recognition of other driving vehicles and the recognition of the lane the sensing vehicle is on at the moment.

%For instance, vehicles which drive through the cities could sense parking spaces while they drive by. There already exists some work on this subject. Mathur et al. \cite{Mathur:2010:PDS:1814433.1814448} developed their system "ParkNet" which can sense parking spaces using an ultrasonic range finder. While the car is driving through the city, the sensor continuously measures the distance to the nearest obstacle on the right side of the road (in many cases a parking car). Using this information and GPS measurements the parking space counts and the parking occupancy rates can be derived by an accuracy of over 90\%. However, such mobile sensing systems currently only work on single lane roads. To work in real world scenarios, lane detection has to be incorporated. 
%Most current approaches on lane detection use cameras which analyse the road markings in the captured pictures. 



\section{Detailed Approach and Goals}
The overall aim of this thesis is to evaluate if it is possible to reach a sufficient high accuracy in road side parking detection on multi lane roads using a sensing vehicle which drives through the city and senses parking spaces while driving by. 
For the parking detection an optical distance sensor will be used to measure the distance to the nearest obstacle on the right side of the road (in many cases a parking car). This sensor will be mounted on the co-driver's side of the car and will continuously measure the distance while the car is driving. Furthermore, a GPS sensor will be used to include the spatial information of the vehicle. Using these measurements, the prototype should reason about the parking situation at the right side of the road. 

%In a first step, the sensor will be tested on accuracy and reliability. 
In a first step, after the sensors have been mounted on a car, test measurements will be collected while driving through some selected streets in Linz, Austria. The test scenes should include single lane- as well as multi lane streets and measurements in all streets should be done several times. To determine the ground truth of the parking availability a camera will be used, which takes pictures of the parking situation at the street during the tests.

%After these measurements have been taken, first an algorithm should be developed which only works for single lane streets. The algorithm should be evaluated on the test measurements and the result of the evaluation are parking space count accuracy and parking occupancy rate accuracy. This can later be used as baseline in comparison to multi lane road detection. 

After these measurements have been taken, the measurements should be analyzed, pre-processed, and then an algorithm should be developed (or learned) to classify the current parking situation. An important part of the algorithm will be the handling of multi lane roads, because there are many special cases which have to be considered. First of all, the lane in which the sensing vehicle is going has to be detected and has to be incorporated in the algorithm. Furthermore, the system should also detect when the car overtakes another driving car, because this could lead to falsely detected parking spaces.

In a final step, possible approaches should be evaluated, which can further enhance the results. For example, the cooperation of multiple sensing vehicles which are going at the same time at the same street can maybe help to increase the parking detection accuracy. Finally, the results of single- and multi lane detection should be evaluated and compared in terms of parking space count accuracy and parking occupancy rate accuracy. 
%Furthermore, the lane detection algorithm based on the distance sensor will be evaluated on its performance.

%\newpage

%\section{Literature}

%In this thesis, the parking situation should be sensed differently. Vehicles should be used as mobile sensors, and while they drive through the city, the collected sensor measurements should be used to derive the current parking availability situation. Such a system should have the advantage to be more cost effective and should also provide a sufficient accurate estimate of the current parking availability in real time.

\newpage

\section{Milestones}

%Table \ref{tab:milestones} shows the planned milestones of this project.

\begin{table}[h]
\centering
\begin{tabular}{r|l}
\emph{Date} & \emph{Milestone} \\\hline
21.04.2017 & Hardware is available \\
05.05.2017 & Hardware parts work together and sensor data can be retrieved. \\
 & Tests regarding the accuracy and range of the optical distance sensor \\
 & are taken. \\
12.05.2017 & The sensors are mounted on the car. \\
19.05.2017 & Test data has been collected on a single lane road. \\
31.06.2017 & Single lane parking detection is implemented and has been evaluated. \\
15.07.2016 & Test data has been collected on multi lane roads in different scenarios. \\
01.09.2016 & Lane detection algorithm has been implemented and evaluated. \\
31.10.2017 & Parking detection on multi lane road has been implemented and \\
 & evaluated \\
30.11.2017 & Further improvements have been implemented and evaluated \\
31.01.2018 & Submission of the thesis
\end{tabular}
\caption{\label{tab:milestones} Milestones}
\end{table}

%%% appendix
%\appendix

%\chapter{Machine Learning Results}
\label{appendix:ml_results}
\markright{Machine Learning Experiment Results}

This appendix lists the results of all tested traditional machine learning models. The confusion matrices for the random forest classifier (Table \ref{table:best_clf_confusion_matrix_filtered}) and for the two-staged classifier which uses its surroundings (Table \ref{table:surrounding_classifier_confusion_matrix}) are already shown in Chapter \ref{chap:evaluation}. The confusion matrices of the best configuration of all other traditional machine learning models can be found in the following.

%The confusion matrices of the best configuration of each model is provided for the full- and filtered dataset. 
%This appendix lists the results of the best configuration of all discussed traditional machine learning models in the form of confusion matrices. Furthermore, the confusion matrix of the best two-stage classifier which uses surrounding segments is also included.







% decision tree
\begin{table}[h]

\resizebox{\textwidth}{!}{%
\centering
\bgroup
\def\arraystretch{1.4}
\begin{tabular}{| r || c | c | c | c |}
\hline
   Predicted class $\rightarrow$ &
   \textbf{Free Space} & 
   \textbf{Parking Car} &
   \textbf{Overtaking} &
   \textbf{Other Parking} \\
   True class $\downarrow$ &
	 & 
	 &
	 \textbf{Situation} &
	 \textbf{Vehicle} \\
\hline
  \textbf{Free Space} & 11616 & 171 & 15 & 9 \\
\hline
  \textbf{Parking Car} & 283 & 1882 & 50 & 1 \\
\hline
  \textbf{Overtaking Situation} & 59 & 83 & 65 & 0 \\
\hline
  \textbf{Other Parking Vehicle} & 49 & 16 & 0 & 4 \\
\hline

\end{tabular}
\egroup
}

\begin{center}
(a) full dataset\\
\end{center}

\resizebox{\textwidth}{!}{%
\centering
\bgroup
\def\arraystretch{1.4}
\begin{tabular}{| r || c | c | c | c |}
\hline
   Predicted class $\rightarrow$ &
   \textbf{Free Space} & 
   \textbf{Parking Car} &
   \textbf{Overtaking} &
   \textbf{Other Parking} \\
   True class $\downarrow$ &
	 & 
	 &
	 \textbf{Situation} &
	 \textbf{Vehicle} \\
\hline
  \textbf{Free Space} & 6541 & 125 & 1 & 13 \\
\hline
  \textbf{Parking Car} & 200 & 1962 & 13 & 8 \\
\hline
  \textbf{Overtaking Situation} & 14 & 37 & 15 & 0 \\
\hline
  \textbf{Other Parking Vehicle} & 32 & 19 & 0 & 9 \\
\hline

\end{tabular}
\egroup
}

\begin{center}
(b) filtered dataset\\
\end{center}

\caption{Confusion matrices of a \textbf{\textit{decision tree}} (using gini impurity as criterion for node-splitting) applied on (a) the full dataset and (b) the filtered dataset.}
\label{table:cf_tree}
\end{table}






% kNN
\begin{table}

\resizebox{\textwidth}{!}{%
\centering
\bgroup
\def\arraystretch{1.4}
\begin{tabular}{| r || c | c | c | c |}
\hline
   Predicted class $\rightarrow$ &
   \textbf{Free Space} & 
   \textbf{Parking Car} &
   \textbf{Overtaking} &
   \textbf{Other Parking} \\
   True class $\downarrow$ &
	 & 
	 &
	 \textbf{Situation} &
	 \textbf{Vehicle} \\
\hline
  \textbf{Free Space} & 11581 & 231 & 14 & 3 \\
\hline
  \textbf{Parking Car} & 276 & 1915 & 24 & 1 \\
\hline
  \textbf{Overtaking Situation} & 97 & 80 & 30 & 0 \\
\hline
  \textbf{Other Parking Vehicle} & 51 & 14 & 0 & 4 \\
\hline

\end{tabular}
\egroup
}

\begin{center}
(a) full dataset\\
\end{center}

\resizebox{\textwidth}{!}{%
\centering
\bgroup
\def\arraystretch{1.4}
\begin{tabular}{| r || c | c | c | c |}
\hline
   Predicted class $\rightarrow$ &
   \textbf{Free Space} & 
   \textbf{Parking Car} &
   \textbf{Overtaking} &
   \textbf{Other Parking} \\
   True class $\downarrow$ &
	 & 
	 &
	 \textbf{Situation} &
	 \textbf{Vehicle} \\
\hline
  \textbf{Free Space} & 6491 & 181 & 1 & 7 \\
\hline
  \textbf{Parking Car} & 177 & 2003 & 3 & 0 \\
\hline
  \textbf{Overtaking Situation} & 28 & 37 & 1 & 0 \\
\hline
  \textbf{Other Parking Vehicle} & 35 & 19 & 0 & 6 \\
\hline

\end{tabular}
\egroup
}

\begin{center}
(b) filtered dataset\\
\end{center}

\caption{Confusion matrices of a \textbf{\textit{kNN classifier}} (using the five nearest neighbours) applied on (a) the full dataset and (b) the filtered dataset.}
\label{table:cf_knn}
\end{table}





% neural network
\begin{table}

\resizebox{\textwidth}{!}{%
\centering
\bgroup
\def\arraystretch{1.4}
\begin{tabular}{| r || c | c | c | c |}
\hline
   Predicted class $\rightarrow$ &
   \textbf{Free Space} & 
   \textbf{Parking Car} &
   \textbf{Overtaking} &
   \textbf{Other Parking} \\
   True class $\downarrow$ &
	 & 
	 &
	 \textbf{Situation} &
	 \textbf{Vehicle} \\
\hline
  \textbf{Free Space} & 11587 & 214 & 9 & 1 \\
\hline
  \textbf{Parking Car} & 319 & 1879 & 17 & 1 \\
\hline
  \textbf{Overtaking Situation} & 78 & 92 & 37 & 0 \\
\hline
  \textbf{Other Parking Vehicle} & 51 & 18 & 0 & 0 \\
\hline

\end{tabular}
\egroup
}

\begin{center}
(a) full dataset\\
\end{center}

\resizebox{\textwidth}{!}{%
\centering
\bgroup
\def\arraystretch{1.4}
\begin{tabular}{| r || c | c | c | c |}
\hline
   Predicted class $\rightarrow$ &
   \textbf{Free Space} & 
   \textbf{Parking Car} &
   \textbf{Overtaking} &
   \textbf{Other Parking} \\
   True class $\downarrow$ &
	 & 
	 &
	 \textbf{Situation} &
	 \textbf{Vehicle} \\
\hline
  \textbf{Free Space} & 6493 & 181 & 6 & 0 \\
\hline
  \textbf{Parking Car} & 214 & 1966 & 3 & 0 \\
\hline
  \textbf{Overtaking Situation} & 21 & 40 & 5 & 0 \\
\hline
  \textbf{Other Parking Vehicle} & 36 & 24 & 0 & 0 \\
\hline

\end{tabular}
\egroup
}

\begin{center}
(b) filtered dataset\\
\end{center}

\caption{Confusion matrices of a \textbf{\textit{neural network}} (using five layers with each 50 hidden units, trained during 1.000.000 epochs) applied on (a) the full dataset and (b) the filtered dataset.}
\label{table:cb_nn}
\end{table}



% svm
\begin{table}

\resizebox{\textwidth}{!}{%
\centering
\bgroup
\def\arraystretch{1.4}
\begin{tabular}{| r || c | c | c | c |}
\hline
   Predicted class $\rightarrow$ &
   \textbf{Free Space} & 
   \textbf{Parking Car} &
   \textbf{Overtaking} &
   \textbf{Other Parking} \\
   True class $\downarrow$ &
	 & 
	 &
	 \textbf{Situation} &
	 \textbf{Vehicle} \\
\hline
  \textbf{Free Space} & 11717 & 93 & 1 & 0 \\
\hline
  \textbf{Parking Car} & 453 & 1762 & 1 & 0 \\
\hline
  \textbf{Overtaking Situation} & 172 & 35 & 0 & 0 \\
\hline
  \textbf{Other Parking Vehicle} & 56 & 13 & 0 & 0 \\
\hline

\end{tabular}
\egroup
}

\begin{center}
(a) full dataset\\
\end{center}

\resizebox{\textwidth}{!}{%
\centering
\bgroup
\def\arraystretch{1.4}
\begin{tabular}{| r || c | c | c | c |}
\hline
   Predicted class $\rightarrow$ &
   \textbf{Free Space} & 
   \textbf{Parking Car} &
   \textbf{Overtaking} &
   \textbf{Other Parking} \\
   True class $\downarrow$ &
	 & 
	 &
	 \textbf{Situation} &
	 \textbf{Vehicle} \\
\hline
  \textbf{Free Space} & 6607 & 73 & 0 & 0 \\
\hline
  \textbf{Parking Car} & 377 & 1806 & 0 & 0 \\
\hline
  \textbf{Overtaking Situation} & 52 & 14 & 0 & 0 \\
\hline
  \textbf{Other Parking Vehicle} & 42 & 18 & 0 & 0 \\
\hline

\end{tabular}
\egroup
}

\begin{center}
(b) filtered dataset\\
\end{center}

\caption{Confusion matrices of a \textbf{\textit{support vector machine}} applied on (a) the full dataset and (b) the filtered dataset.}
\label{table:cf_svm}
\end{table}




% naive bayes
\begin{table}

\resizebox{\textwidth}{!}{%
\centering
\bgroup
\def\arraystretch{1.4}
\begin{tabular}{| r || c | c | c | c |}
\hline
   Predicted class $\rightarrow$ &
   \textbf{Free Space} & 
   \textbf{Parking Car} &
   \textbf{Overtaking} &
   \textbf{Other Parking} \\
   True class $\downarrow$ &
	 & 
	 &
	 \textbf{Situation} &
	 \textbf{Vehicle} \\
\hline
  \textbf{Free Space} & 5013 & 2159 & 561 & 4078 \\
\hline
  \textbf{Parking Car} & 55 & 1942 & 78 & 141 \\
\hline
  \textbf{Overtaking Situation} & 9 & 89 & 79 & 30 \\
\hline
  \textbf{Other Parking Vehicle} & 1 & 11 & 0 & 57 \\
\hline

\end{tabular}
\egroup
}

\begin{center}
(a) full dataset\\
\end{center}

\resizebox{\textwidth}{!}{%
\centering
\bgroup
\def\arraystretch{1.4}
\begin{tabular}{| r || c | c | c | c |}
\hline
   Predicted class $\rightarrow$ &
   \textbf{Free Space} & 
   \textbf{Parking Car} &
   \textbf{Overtaking} &
   \textbf{Other Parking} \\
   True class $\downarrow$ &
	 & 
	 &
	 \textbf{Situation} &
	 \textbf{Vehicle} \\
\hline
  \textbf{Free Space} & 3288 & 1298 & 147 & 1947 \\
\hline
  \textbf{Parking Car} & 64 & 1926 & 56 & 137 \\
\hline
  \textbf{Overtaking Situation} & 0 & 22 & 29 & 15 \\
\hline
  \textbf{Other Parking Vehicle} & 2 & 10 & 0 & 48 \\
\hline

\end{tabular}
\egroup
}

\begin{center}
(b) filtered dataset\\
\end{center}

\caption{Confusion matrices of a \textbf{\textit{naive bayes classifier}} applied on (a) the full dataset and (b) the filtered dataset.}
\label{table:cf_nb}
\end{table}











%
%
%
%% deep learning neural network
%\begin{table}
%
%\resizebox{\textwidth}{!}{%
%\centering
%\bgroup
%\def\arraystretch{1.4}
%\begin{tabular}{| r || c | c | c | c |}
%\hline
%   Predicted class $\rightarrow$ &
%   \textbf{Free Space} & 
%   \textbf{Parking Car} &
%   \textbf{Overtaking} &
%   \textbf{Other Parking} \\
%   True class $\downarrow$ &
%	 & 
%	 &
%	 \textbf{Situation} &
%	 \textbf{Vehicle} \\
%\hline
%  \textbf{Free Space} & 5013 & 2159 & 561 & 4078 \\
%\hline
%  \textbf{Parking Car} & 55 & 1942 & 78 & 141 \\
%\hline
%  \textbf{Overtaking Situation} & 9 & 89 & 79 & 30 \\
%\hline
%  \textbf{Other Parking Vehicle} & 1 & 11 & 0 & 57 \\
%\hline
%
%\end{tabular}
%\egroup
%}
%
%\begin{center}
%(a) full dataset\\
%\end{center}
%
%\resizebox{\textwidth}{!}{%
%\centering
%\bgroup
%\def\arraystretch{1.4}
%\begin{tabular}{| r || c | c | c | c |}
%\hline
%   Predicted class $\rightarrow$ &
%   \textbf{Free Space} & 
%   \textbf{Parking Car} &
%   \textbf{Overtaking} &
%   \textbf{Other Parking} \\
%   True class $\downarrow$ &
%	 & 
%	 &
%	 \textbf{Situation} &
%	 \textbf{Vehicle} \\
%\hline
%  \textbf{Free Space} & 3288 & 1298 & 147 & 1947 \\
%\hline
%  \textbf{Parking Car} & 64 & 1926 & 56 & 137 \\
%\hline
%  \textbf{Overtaking Situation} & 0 & 22 & 29 & 15 \\
%\hline
%  \textbf{Other Parking Vehicle} & 2 & 10 & 0 & 48 \\
%\hline
%
%\end{tabular}
%\egroup
%}
%
%\begin{center}
%(b) filtered dataset\\
%\end{center}
%
%\caption{Confusion matrices of a \textbf{\textit{naive bayes classifier}} applied on (a) the full dataset and (b) the filtered dataset.}
%\label{table:cf_nb}
%\end{table}

% \chapter*{Abbreviations}
\markright{Abbreviations}

\begin{description}
\setlength{\itemsep}{-11pt}
\setlength{\leftmargin}{900pt}

%C
\item[CNN] Convolutional Neural Network

%D
\item[DL] Deep Learning

%G
\item[GNNS] Global Navigation Satellite System
\item[GPIO] General Purpose Input Output
\item[GPS] Global Positioning System

%I
\item[I2C] Inter Integrated Circuit

%L
\item[LIDAR] Light detection and ranging

\end{description}


\bibliographystyle{plain}
\bibliography{literature}


% curriculum vitae
%\input{structure/cv}
\end{document}
